\section{The \acr{LLVM} \acr{IR} Project}
\acr{LLVM} began as an open-source compiler infrastructure project at the University of Illinois in 2000, originally designed to provide a modern, \acr{SSA}-based compilation technique for arbitrary static and dynamic programming languages.

The design and core concepts developed in the \acr{LLVM} research project have proven highly advantageous for language implementation and compilation due to its unified \acr{ISA} and the \textbf{conceptual similarity} between \acr{LLVM} \acr{IR} and real assembly languages. Consequently, the project has evolved into a comprehensive ecosystem of distinct subprojects providing reusable libraries and \acr{ISA} specifications.

\subsection{The Compilation Pipeline}
Traditional compilation flows from source code $\rightarrow$ \acr{AST} $\rightarrow$ assembly $\rightarrow$ machine code. Given the architecture diversity discussed in the previous section, this creates the $M \times N$ problem: $M$ languages each needing $N$ architecture-specific backends.

To address this challenge, \acr{LLVM} introduced a \textbf{unified abstract low-level language} with a common \acr{ISA} that is structurally analogous to native assembly languages. The compilation process (translation to machine code for each architecture) is then executed by the \acr{LLVM} compiler infrastructure, which acts as a virtual machine for this abstract representation.

\acr{LLVM} \acr{IR} is therefore a low-level intermediate language that resembles assembly but with several critical differences. First, it is \textbf{\textit{strongly typed}}---every value has an explicit type (\texttt{i32}, \texttt{float}, \texttt{ptr}, etc.), enabling compile-time type checking and optimization opportunities. Second, it is \textbf{\textit{platform-independent}}, meaning the same \acr{IR} code works for \acr{x86}, \acr{ARM}, \acr{GPU}, and \acr{WASM} targets without modification. Third, it enforces \textbf{\textit{\acr{SSA} form}}, where every variable is assigned exactly once, enabling aggressive optimizations by simplifying data flow analysis. Fourth, it provides \textbf{\textit{infinite registers}}---virtual registers prefixed with \texttt{\%} that the backend maps to physical hardware registers during code generation. Finally, it uses \textbf{\textit{explicit control flow}} through structured basic blocks and explicit branch instructions instead of unstructured goto statements.

Example \acr{LLVM} \acr{IR} function:
\begin{lstlisting}[style=ir]
define i32 @add(i32 %a, i32 %b) {
entry:
  %result = add i32 %a, %b
  ret i32 %result
}
\end{lstlisting}

This representation is simultaneously human-readable (unlike binary machine code), optimizable (unlike high-level source code), and portable (unlike assembly language).

\subsection{Optimization Pipeline}
As mentioned above, the \acr{LLVM} infrastructure compiles (translates to machine code) the intermediate representation of instructions for each specific architecture. However, thanks to the use of the \acr{SSA} concept, the resulting instruction graph structure can be optimized—meaning code with better performance characteristics can be generated for the target architecture. Unfortunately, there is no simple universal definition of \textit{optimal code} with respect to performance; optimality is determined by the specific requirements and design goals of each application.

\acr{LLVM} supports six optimization levels, each applying progressively more aggressive transformations:

\begin{description}
  \item[\texttt{-O0} (No optimization)] This level performs \textbf{no optimizations whatsoever}. The compiler generates machine code that directly corresponds to the \acr{IR} instructions, making the relationship between source code and assembly trivial to understand. This is essential for debugging, as it preserves variable locations, call stack integrity, and stepping behavior. Compilation is fastest at this level. Use \texttt{-O0} during active development and debugging sessions.

  \item[\texttt{-O1} (Basic optimization)] Applies \textbf{fundamental optimizations} that improve performance with minimal compilation time overhead. These include \textbf{constant folding} (evaluating constant expressions at compile time, e.g., \texttt{x = 3 * 4} becomes \texttt{x = 12}), \textbf{dead code elimination} (removing unreachable code and unused variables), \textbf{simple inlining}\footnote{In compiler terminology, \textit{inlining} refers to replacing a function call with the actual body of the function at the call site, eliminating call overhead and enabling further optimizations. This differs from the CSS concept of ``inline'' (inline styles or inline elements). For example, if \texttt{add(a, b)} calls a function that returns \texttt{a + b}, inlining replaces the call with the direct computation \texttt{a + b}, removing the function call entirely.} (inlining trivial functions such as getters, setters, and small wrappers), and \textbf{expression simplification} (reducing algebraic expressions, e.g., \texttt{x * 1} becomes \texttt{x}). Use \texttt{-O1} when you want some performance improvement while maintaining reasonable compile times and debuggability.

  \item[\texttt{-O2} (Moderate optimization)] The \textbf{recommended level for production code}. Applies all \texttt{-O1} optimizations plus \textbf{loop optimizations} (loop unrolling, loop invariant code motion, strength reduction), \textbf{vectorization} (automatically converting scalar operations to \acr{SIMD} instructions such as \acr{SSE} and \acr{AVX} on \acr{x86} or \acr{NEON} on \acr{ARM}), \textbf{function inlining} (aggressive inlining of hot functions based on profitability heuristics), \textbf{\acr{CSE}} (reusing results of repeated calculations), \textbf{register allocation optimization} (better use of physical registers, reducing memory spills), and \textbf{instruction scheduling} (reordering instructions to minimize \acr{CPU} pipeline stalls). This level provides an excellent balance between compilation time and runtime performance. Most production software uses \texttt{-O2}.

  \item[\texttt{-O3} (Aggressive optimization)] Applies all \texttt{-O2} optimizations plus \textbf{aggressive transformations that may increase code size}, including \textbf{aggressive loop transformations} (loop interchange, loop fusion, polyhedral optimization), \textbf{aggressive inlining} (inlining larger functions, potentially duplicating code across call sites), \textbf{\acr{IPO}} (analyzing and optimizing across function boundaries), \textbf{tail call optimization} (converting recursive calls to iterative loops where possible), and \textbf{speculative execution} (duplicating code paths to enable better scheduling). Use \texttt{-O3} for compute-intensive workloads (numerical computing, machine learning, \acr{GPU} kernels) where performance is critical and binary size is not a concern. Note that \texttt{-O3} can occasionally introduce numerical instability in floating-point code due to aggressive reassociation.

  \item[\texttt{-Os} (Optimize for size)] Applies \textbf{optimizations that reduce binary size} while maintaining reasonable performance. This level enables most \texttt{-O2} optimizations except those that increase code size, disables aggressive inlining and loop unrolling, and prefers smaller instruction encodings when semantically equivalent. Use \texttt{-Os} for embedded systems with limited flash/ROM, \acr{WASM} modules (smaller downloads), or when deploying to size-constrained environments (Docker containers, mobile apps).

  \item[\texttt{-Oz} (Aggressively optimize for size)] Provides even more \textbf{aggressive size reduction than \texttt{-Os}} by disabling vectorization (\acr{SIMD} instructions are larger than scalar equivalents), applying aggressive function outlining (extracting common code sequences into shared functions), and potentially sacrificing performance for size reduction. Use \texttt{-Oz} when binary size is the absolute priority, such as bootloaders, firmware, or extremely resource-constrained \acr{IoT} devices.
\end{description}

In sections~\ref{sec:llvm-optimization-flags} and~\ref{sec:gpu-compilation}, we comprehensively demonstrate how to apply each of these optimization strategies to \acr{CPU} and \acr{GPU} code generation workflows.

\subsection{Portability}
One of \acr{LLVM} \acr{IR}'s most significant advantages is \textbf{true cross-platform portability}. Unlike traditional assembly languages where \texttt{x86} code, \acr{ARM} code, and \acr{RISC-V} code are fundamentally incompatible, \acr{LLVM} \acr{IR} provides a \textbf{write once, compile anywhere} paradigm.

The same \acr{IR} code can be compiled to native machine code for \acr{x86-64}, \acr{ARM}, \acr{RISC-V}, \acr{MIPS}, PowerPC, \acr{GPU} architectures (NVIDIA, AMD, Intel), and even \acr{WASM} targets \textit{without modification}. The \acr{LLVM} backend automatically handles all architecture-specific details—instruction selection, register allocation, calling conventions, and alignment requirements—transparently during code generation.

This portability eliminates the traditional problem where developers must maintain separate assembly codebases for each target platform. Instead of writing and debugging architecture-specific implementations, you write \acr{IR} once, and the \acr{LLVM} infrastructure ensures correct compilation for all supported targets. This dramatically reduces maintenance burden and enables seamless cross-platform deployment.

For performance-critical sections traditionally written in inline assembly (cryptography routines, vectorized numeric kernels, system-level operations), \acr{LLVM} \acr{IR} offers comparable performance while maintaining portability. The \acr{IR} abstracts away hardware differences while preserving low-level control over memory layout, instruction sequencing, and optimization hints.

As programming technologies evolve, the \acr{LLVM} project continuously adapts by introducing \textbf{specialized \acr{IR} dialects} and extensions that support emerging compilation requirements. For instance, \textbf{\acr{NVVM} \acr{IR}} (NVIDIA Virtual Machine) extends standard \acr{LLVM} \acr{IR} with \acr{GPU}-specific intrinsics for thread indexing, shared memory management, and synchronization primitives required for \acr{CUDA} compilation. Similarly, \textbf{\acr{AMDGPU} \acr{IR}} provides \acr{AMD}-specific extensions, and \textbf{\acr{SPIR-V}} serves as a portable intermediate representation for \acr{OpenCL} and Vulkan compute shaders. These dialects maintain full compatibility with the core \acr{LLVM} infrastructure while enabling platform-specific optimizations and features, ensuring that \acr{LLVM} remains current with evolving hardware capabilities and programming paradigms without sacrificing its foundational portability guarantees.
